\begin{homeworkProblem}[2]
  Construa o conjunto de Cantor $\cali{C}$ da seguinte forma: começamos com $C_0=[0,1]$, removemos seu terço médio para obter $C_1=\left[0,\frac{1}{3}\right]\cup\left[\frac{2}{3},1\right]$, e, indutivamente, $C_k$ é obtido retirando o terço médio de cada intervalo que compõe $C_{k-1}$ de forma a obter $2^{k+1}$ intervalos fechados de comprimento $\frac{1}{3^k}$. Definimos então $\cali{C}=\bigcap_{n=1}^{\infty}C_n$. Mostre as seguintes propriedades
  \begin{enumerate}[a)]
    \item $\cali{C}$ é fechado, tem medida $m(\cali{C})=0$ e não possui pontos insolados.
    \item Um ponto $x\in[0,1]$ pertenece a $\cali{C}$ se, somente se, sua expressão na base $3$ possui apenas os dígitos $0$ e $2$, ou seja, existe uma sequência $\{d_n\}_{n\in\mathbb{Z}^{+}}$ tal que $d_n\in\{0,2\}$ para todo $n\in\mathbb{Z}^{+}$ e $x=\sum_{n\in\mathbb{Z}^{+}}\frac{d_n}{3^n}$.
  \end{enumerate}
  \begin{solution}
    \begin{enumerate}[a)]
      \item Note que, para cada $n\in\mathbb{Z}^{+}$, $C_n$ é uma união finita de intervalos fechados. Portanto, $C_n$ é fechado. Logo, como interseção arbitrária de conjuntos fechados é fechada, temos que $\cali{C}\bigcap_{n=1}^{\infty}C_n$ é fechado.\\
        Além disso, $C_n$ é formado por $2^n$ intervalos fechados, cada um de comprimento $1/3^n$. Logo,
        \begin{align*}
          m(C_n)=2^n\frac{1}{3^n}=\left(\frac{2}{3}\right)^n.
        \end{align*}
        Portanto, como $\cali{C}\subseteq C_n$ para todo $n\in\mathbb{Z}^{+}$, temos que
        \begin{align*}
          m(\cali{C})\leq m(C_n)=\left(\frac{2}{3}\right)^n,\quad\text{para todo }n\in\mathbb{Z}^{+}.
        \end{align*}
        Isto é, que $m(\cali{C})=0$.\\
        Agora, mostraremos que $C$ não possui pontos isolados.\\
        Seja $x\in C$. Note que para cada $n\in\mathbb{Z}^{+}$, como $x\in C_n$, então existe um único intervalo $I_n$ que compõe $C_n$ tal que $x\in I_n$ e $|I_n|=1/3^n$.\\
        Logo, como $x\in C_{n+1}$, temos que $x$ pertence a um dos dois intervalos obtidos ao retirar o terço médio de $I_n$, definamos $J_n$ o outro intervalo.\\
        Observe que, um dos extremos de $J_n$ pertence a $\cali{C}$, pois os extremos dos intervalos da construção de Cantor permanecem em todo $C_n$. Escolhemos esse extremo e o denotamos por $y_n$.\\
        Portanto, temos que $y_n\in C$, $y_n\neq x$ e
        \begin{align*}
          |x-y_n|\leq |I_n|=\frac{1}{3^n}.
        \end{align*}
        O que nos permite garantir que o $\lim_{n\to\infty}y_n=x$. Assim, como $x\in \cali{C}$ é arbitrário, temos que todo ponto de $\cali{C}$ é ponto de acumulação de $\cali{C}$, isto é, $\cali{C}$ não possui pontos isolados.
      \item Suponha inicialmente que $x\in \cali{C}$. Logo, note que $x\in \left[ 0,\frac{1}{3} \right]$ ou $x\in\left[ \frac{2}{3},1 \right]$, Logo, pela construção do conjunto, temos que $d_1$ está determinado pelos extremos dos intervalos o que nos leva a considerar o seguinte
        \begin{itemize}
          \item Se $x=1$, então a expressão na base $3$ de $x$ é $2.\overline{2}$.
          \item Se $x=\frac{1}{3}$, então a expressão na base $3$ de $x$ é $0.\overline{2}$.
          \item No outro caso, temos que $d_1\in\{0,2\}$. 
        \end{itemize}
        Assim, raciocinando indutivamente, temos que para cada $n\in\mathbb{Z}^{+}$, o coeficiente $d_n\in\{0,2\}$.\\
        Reciprocamente, suponha que existe uma sequência $\{d_n\}_{n\in\mathbb{Z}^{+}}$ tal que $d_n\in{0,2}$ para todo $n\in\mathbb{Z}^{+}$.\\
        Como vimos no raciocínio anterior, temos que cada coeficiente $d_n$ está indicando o extremo do intervalo de $C_n$ no que $x$ pertence, pelo que podemos garantir que $x\in C_n$ para todo $n\in\mathbb{Z}^{+}$, isto é, que $x\in\cali{C}$.
    \end{enumerate}
  \end{solution}
\end{homeworkProblem}
